A entrega frequente de produtos e sistemas digitais envolve a coordenação
de processos técnicos e organizacionais, desde o alinhamento de requisitos
até a preparação da infraestrutura. A busca por entregas que gerem valor
e mantenham a qualidade tem motivado o uso de abordagens ágeis e integradas,
como alternativas às limitações de modelos sequenciais em determinados
contextos. Amaro, Pereira e Mira da Silva~\citeano{amaro2025} discutem a
articulação das capacidades DevOps ao longo do ciclo de vida do software.

DevOps compreende práticas destinadas a integrar e automatizar atividades
das equipes de desenvolvimento de software e de operações de TI. Seu
objetivo é favorecer entregas contínuas, conciliando agilidade, qualidade,
estabilidade e segurança. Além da adoção de ferramentas, envolve uma
cultura organizacional orientada à colaboração e ao compartilhamento de
responsabilidades \cite{azad2023,amaro2026}.

Uma das barreiras historicamente associadas à adoção de DevOps é a separação
de responsabilidades entre desenvolvimento e operações em equipes isoladas,
frequentemente denominadas \emph{silos}. Enquanto o desenvolvimento prioriza a
disponibilização de novas funcionalidades, as operações precisam assegurar
a estabilidade do ambiente de produção. A falta de coordenação entre essas
atividades pode gerar conflitos de prioridades e dificultar as entregas.
Azad e Hyrynsalmi~\citeano{azad2023} destacam a integração entre
desenvolvimento e operações como fator relevante para o desempenho dos
processos de engenharia de software.

Para reduzir essas barreiras, a cultura DevOps articula automação,
integração contínua (CI) e entrega contínua (CD) em um fluxo recorrente
de trabalho. Essas práticas não se restringem a uma implementação isolada:
sua continuidade depende da integração de capacidades, processos e
mecanismos de feedback ao longo do ciclo de vida do software
\cite{amaro2025,amaro2026}. À medida que as entregas se tornam mais
frequentes, acompanhar o comportamento dos sistemas em produção também
ganha importância para orientar as decisões de engenharia.
